\begin{proposition}
\label{proposition-detecting-bounded-above}
Let $X$ be a quasi-compact and quasi-separated scheme. Let
$G \in D_{perf}(\mathcal{O}_X)$ be a perfect complex which generates 
$D_\QCoh (\mathcal{O}_X)$. Let $E \in D_\QCoh (\mathcal{O}_X)$.
The following are equivalent
\begin{enumerate}
\item $E \in D^-_\QCoh (\mathcal{O}_X)$,
\item $\Hom_{D(\mathcal{O}_X)}(G[-i], E) = 0$ for $i \gg 0$,
\item $\Ext^i_X(G, E) = 0$ for $i \gg 0$,
\item $R\Hom_X(G, E)$ is in $D^-(\mathbf{Z})$,
\item $H^i(X, G^\vee \otimes_{\mathcal{O}_X}^\mathbf{L} E) = 0$
for $i \gg 0$,
\item $R\Gamma(X, G^\vee \otimes_{\mathcal{O}_X}^\mathbf{L} E)$
is in $D^-(\mathbf{Z})$,
\item for every perfect object $P$ of $D(\mathcal{O}_X)$
\begin{enumerate}
\item the assertions (2), (3), (4) hold with $G$ replaced by $P$, and
\item $H^i(X, P \otimes_{\mathcal{O}_X}^\mathbf{L} E) = 0$ for $i \gg 0$,
\item $R\Gamma(X, P \otimes_{\mathcal{O}_X}^\mathbf{L} E)$
is in $D^-(\mathbf{Z})$.
\end{enumerate}
\end{enumerate}
\end{proposition}

\begin{proof}
Assume (1). Since
$\Hom_{D(\mathcal{O}_X)}(G[-i], E) = \Hom_{D(\mathcal{O}_X)}(G, E[i])$
we see that this is zero for $i \gg 0$ by
Lemma \ref{lemma-ext-from-perfect-into-bounded-QCoh}. This proves
that (1) implies (2).

\medskip\noindent
Parts (2), (3), (4) are equivalent by the discussion in
Cohomology, Section \ref{cohomology-section-global-RHom}.
Part (5) and (6) are equivalent as $H^i(X, -) = H^i(R\Gamma(X, -))$
by definition. The equivalent conditions (2), (3), (4) are
equivalent to the equivalent conditions (5), (6) by
Cohomology, Lemma \ref{cohomology-lemma-dual-perfect-complex}
and the fact that $(G[-i])^\vee = G^\vee[i]$.

\medskip\noindent
It is clear that (7) implies (2). Conversely, 
let us prove that the equivalent conditions (2) -- (6) imply (7).
Recall that $G$ is a  classical generator for $D_{perf}(\mathcal{O}_X)$ by
Remark \ref{remark-classical-generator}.
For $P \in D_{perf}(\mathcal{O}_X)$ let $T(P)$ be the assertion that
$R\Hom_X(P, E)$ is in $D^-(\mathbf{Z})$.
Clearly, $T$ is inherited by direct sums,
satisfies the 2-out-of-three property for distinguished
triangles, is inherited by direct summands, and is preserved by shifts.
Hence by Derived Categories, Remark \ref{derived-remark-check-on-generator}
we see that (4) implies $T$ holds on all of $D_{perf}(\mathcal{O}_X)$.
The same argument works for all other properties, except that for property
(7)(b) and (7)(c) we also use that $P \mapsto P^\vee$ is a self
equivalence of $D_{perf}(\mathcal{O}_X)$. Small detail omitted.

\medskip\noindent
We will prove the equivalent conditions (2) -- (7) imply (1)
using the induction principle of
Cohomology of Schemes, Lemma \ref{coherent-lemma-induction-principle}.

\medskip\noindent
First, we prove (2) -- (7) $\Rightarrow$ (1) if $X$ is affine.
Set $P = \mathcal{O}_X[0]$. From (7) we obtain $H^i (X, E) = 0$ for $i \gg 0$.
Hence (1) follows since $E$ is determined by $R\Gamma (X, E)$,
see Lemma \ref{lemma-affine-compare-bounded}.

\medskip\noindent
Now assume $X = U \cup V$ with $U$ a quasi-compact open of $X$ and
$V$ an affine open, and assume the implication (2) -- (7) $\Rightarrow$ (1)
is known for the schemes $U$, $V$, and $U \cap V$.
Suppose $E \in D_\QCoh(\mathcal{O}_X)$ satisfies (2) -- (7).
By Lemma \ref{lemma-direct-summand-of-a-restriction} and
Theorem \ref{theorem-bondal-van-den-Bergh} there exists a perfect
complex $Q$ on $X$ such that $Q|_U$ generates $D_\QCoh (\mathcal{O}_U)$. 
Let $f_1, \dots , f_r \in \Gamma (V, \mathcal{O}_V)$
be such that $V \setminus U = V(f_1, \dots , f_r)$ as subsets of
$V$. Let $K \in D_{perf}(\mathcal{O}_V)$ be the object
corresponding to the Koszul complex on $f_1, \dots , f_r$.
Let $K' \in D_{perf}(\mathcal{O}_X)$ be
\begin{equation}
\label{equation-detecting-bounded-above}
K' = R (V \to X)_* K = R (V \to X)_! K,
\end{equation}
see Cohomology, Lemmas \ref{cohomology-lemma-pushforward-restriction} and
\ref{cohomology-lemma-pushforward-perfect}. This is a perfect
complex on $X$ supported on the closed set $X \setminus U \subset V$
and isomorphic to $K$ on $V$. By assumption, we know
$R\Hom_{\mathcal{O}_X}(Q, E)$ and
$R\Hom_{\mathcal{O}_X}(K', E)$ are bounded above.

\medskip\noindent
By the second description of $K'$ in (\ref{equation-detecting-bounded-above})
we have
$$
\Hom_{D(\mathcal{O}_V)}(K[-i], E|_V) = \Hom_{D(\mathcal{O}_X)}(K'[-i], E) = 0
$$
for $i \gg 0$. Therefore, we may apply
Lemma \ref{lemma-orthogonal-koszul-first-variant} to $E|_V$ to
obtain an integer $a$ such that
$\tau_{\geq a}(E|_V) = \tau_{\geq a} R (U \cap V \to V)_* (E|_{U \cap V})$.
Then $\tau_{\geq a} E = \tau_{\geq a} R (U \to X)_* (E |_U)$
(check that the canonical map is an isomorphism after restricting to
$U$ and to $V$). Hence using Lemma \ref{lemma-bounded-truncation}
twice we see that
$$
\Hom_{D(\mathcal{O}_U)}(Q|_U [-i], E|_U) =
\Hom_{D(\mathcal{O}_X)}(Q[-i], R (U \to X)_* (E|_U)) = 0
$$
for $i \gg 0$. Since the Proposition holds for $U$ and the generator
$Q|_U$, we have $E|_U \in D^-_\QCoh(\mathcal{O}_U)$. But then since
the functor $R (U \to X)_*$ preserves $D^-_\QCoh$ 
(by Lemma \ref{lemma-quasi-coherence-direct-image}), we get
$\tau_{\geq a}E \in D^-_\QCoh(\mathcal{O}_X)$. Thus 
$E \in D^-_\QCoh (\mathcal{O}_X)$. 
\end{proof}